\chapter{حساب التفاضل}\label{ch:b2:diffcalc}

ينضج حساب \omterm{def:b2:diffcalc:differential}{التفاضل} بمتغيرين، الذي رأيناه في مجلد السنة الأولى، فيصير
حساب \omterm{def:b2:diffcalc:differential}{تفاضل} التطبيقات بين \omterm{def:b2:nvs:norm}{الفضاءات المعيارية}: \emph{\omterm{def:b2:diffcalc:differential}{التفاضل}} بوصفه
أفضل تقريب خطي، وقاعدة السلسلة في عموميتها التامة، ومبرهنة تناظر
شوارتز \emph{مبرهنة} هذه المرة، وصيغ تايلور، والتحليل الكامل من
الرتبة الثانية للقيم القصوى. أما مبرهنة الدالة العكسية، تاج النظرية،
فتُعرض مع بيان صريح لاستراتيجية برهانها --- نقطة بناخ الثابتة.

في كل ما يلي، $U$ جزء مفتوح من $\R^n$ (أو من فضاء معياري؛ فالحالة
المنتهية البعد تحمل كل الأفكار)، $f \colon U \to
\R^m$.

\section{التفاضل}

\begin{definition}\label{def:b2:diffcalc:differential}
نقول إن $f$ \emph{قابلة للتفاضل}\index{قابلة للتفاضل} في $a$ إذا وُجد
تطبيق خطي (\omterm{def:b2:metric:continuity}{متصل}) $\dd f_a \colon \R^n \to \R^m$ يحقق
\[
f(a + h) = f(a) + \dd f_a(h) + o(\norm h)
\qquad (h \to 0).
\]
التطبيق $\dd f_a$، وهو \emph{تفاضل}\index{تفاضل}
$f$ في $a$، وحيد؛ ومصفوفته في الأساسين القانونيين هي
\emph{مصفوفة جاكوبي}\index{مصفوفة جاكوبي} $J_f(a) =
\bigl(\frac{\partial f_i}{\partial x_j}(a)\bigr)$.
قابلية التفاضل تستلزم الاتصال ووجود جميع المشتقات
الاتجاهية $\dd f_a(v) = \lim_{t\to0}\frac{f(a + tv)
- f(a)}{t}$؛ والعكس غير صحيح
(\cref{exo:b2:diffcalc:2}). ومن أجل $m = 1$، $\dd f_a(h) = \langle
\nabla f(a), h\rangle$: وهو تدرج السنة الأولى، مفهوما الآن على أنه
المتجهة الممثلة للتفاضل.
\end{definition}

\begin{theorem}[محك $C^1$]\label{thm:b2:diffcalc:c1}
إذا وُجدت جميع المشتقات الجزئية للدالة $f$ على $U$ وكانت متصلة
في $a$، فإن $f$ \omterm{def:b2:diffcalc:differential}{قابلة للتفاضل} في $a$. ومنه فإن ``$C^1$ على $U$'' ---
أي اتصال المشتقات الجزئية --- يستلزم قابلية \omterm{def:b2:diffcalc:differential}{التفاضل} في كل نقطة،
مع \omterm{def:b2:diffcalc:differential}{تفاضل} \omterm{def:b2:metric:continuity}{متصل}.
\end{theorem}

\begin{proof}
مركبة مركبة (يكفي $m = 1$). برهان السنة الأولى في حالة متغيرين ---
نحرك إحداثية واحدة في كل مرة، ونطبق مبرهنة القيمة المتوسطة بمتغير
واحد على كل ضلع، ونستعمل اتصال المشتقات الجزئية
في $a$ --- يتعمم حرفيا إلى $n$ ضلعا:
\[
f(a + h) - f(a) = \sum_{j=1}^{n} \bigl(f(a + h^{(j)}) - f(a +
h^{(j-1)})\bigr)
= \sum_j h_j\,\frac{\partial f}{\partial x_j}(\xi_j),
\]
حيث $h^{(j)}$ يجمد الإحداثيات $j$ الأولى من $h$، و$\xi_j$
يقع على الضلع رقم $j$؛ والاتصال يحول كل
$\frac{\partial f}{\partial x_j}(\xi_j)$ إلى $\frac{\partial
f}{\partial x_j}(a) + o(1)$، فيكون الخطأ $o(\norm h)$.
\end{proof}

\begin{theorem}[قاعدة السلسلة]\label{thm:b2:diffcalc:chainrule}
إذا كانت $f$ \omterm{def:b2:diffcalc:differential}{قابلة للتفاضل} في $a$ و$g$ في $f(a)$، فإن $g \circ f$
\omterm{def:b2:diffcalc:differential}{قابلة للتفاضل} في $a$ ولدينا
\[
\dd(g \circ f)_a = \dd g_{f(a)} \circ \dd f_a ,
\qquad
J_{g\circ f}(a) = J_g\bigl(f(a)\bigr)\,J_f(a) :
\]
أي إن مصفوفات جاكوبي تتضارب.\index{قاعدة السلسلة}
\end{theorem}

\begin{proof}
نكتب $f(a + h) = f(a) + \dd f_a(h) + \norm h\,\varepsilon_1(h)$
و$g(b + k) = g(b) + \dd g_b(k) + \norm k\,\varepsilon_2(k)$ مع
$b = f(a)$، $k = k(h) = \dd f_a(h) + \norm h \varepsilon_1(h)$.
وبالتعويض،
\[
g(f(a+h)) = g(b) + \dd g_b\bigl(\dd f_a(h)\bigr)
+ \norm h\,\dd g_b(\varepsilon_1(h)) + \norm{k}\,\varepsilon_2(k),
\]
وحدّا الخطأ كلاهما $o(\norm h)$: الأول لأن $\dd g_b$
\omterm{def:b2:metric:continuity}{متصل} و$\varepsilon_1 \to 0$؛ والثاني لأن $\norm k
\leq C\norm h$ (تطبيق خطي محدود زائد حد صغير) و$\varepsilon_2(k) \to 0$ عندما $h \to 0$.
\end{proof}

\begin{example}[الدوال القطرية، مرة واحدة وإلى الأبد]\label{ex:b2:diffcalc:radial}
نضع $r(x) = \norm x_2$ على $\R^n\setminus\{0\}$ و$f = g \circ
r$ حيث $g$ دالة $C^1$ بمتغير واحد. أولا، $r$
\omterm{def:b2:diffcalc:differential}{قابلة للتفاضل} بعيدا عن $0$: انطلاقا من $r^2 = \sum x_i^2$،
\[
\frac{\partial r}{\partial x_i} = \frac{x_i}{r},
\qquad\text{أي}\qquad
\nabla r(x) = \frac{x}{\norm x} ,
\]
وهي المتجهة القطرية الواحدية (نفاضل $r^2$ ونقسم --- أو نطبق قاعدة
السلسلة على $\sqrt{\cdot}$). ثم تعطي قاعدة السلسلة، من أجل كل دالة قطرية،
\[
\nabla f(x) = g'\bigl(\norm x\bigr)\,\frac{x}{\norm x} .
\]
وإليك حالة محسوبة: $g(r) = \frac1r$ يعطي $\nabla\frac{1}{\norm
x} = -\frac{x}{\norm x^3}$، وهو حقل مقلوب المربع في
الجاذبية والكهرباء الساكنة --- اتجاهه قطري وقيمته
$\frac{1}{\norm x^2}$. والخلاصة: تدرجات الدوال القطرية قطرية لأن مجموعات المستوى
كرات، ولأن التدرج عمودي على مجموعات المستوى؛ أما في $x = 0$ فعلى
النقيض، $r$ \emph{ليست} \omterm{def:b2:diffcalc:differential}{قابلة للتفاضل} (لا يوافق أي تطبيق خطي
مرشح القيمة $\norm h$ من جميع الاتجاهات) --- فالمنحنيات القطرية الملساء
تحتاج إلى $g'(0) = 0$ كي تعبر المبدأ بسلاسة.
\end{example}

\begin{theorem}[متراجحة القيمة المتوسطة]\label{thm:b2:diffcalc:mvi}
لتكن $f$ \omterm{def:b2:diffcalc:differential}{قابلة للتفاضل} على $U$، وليكن القطعة $\intcc{a}{b}
= \{a + t(b-a)\}$ محتواة في $U$. عندئذ
\[
\norm{f(b) - f(a)} \leq \norm{b - a}\,
\sup_{x \in \intcc{a}{b}} \vertiii{\dd f_x} .
\]
وبوجه خاص، كل تطبيق قابل \omterm{def:b2:diffcalc:differential}{للتفاضل} تفاضله معدوم على مفتوح
\emph{مترابط} يكون ثابتا.
\end{theorem}

\begin{proof}
الدالة $\varphi(t) = f(a + t(b-a))$ \omterm{def:b2:diffcalc:differential}{قابلة للتفاضل} على
$\intcc{0}{1}$ ولدينا $\varphi'(t) = \dd f_{a + t(b-a)}(b - a)$
(قاعدة السلسلة)، ومعيارها $\leq M\norm{b-a}$ حيث $M$ هو الحد الأعلى المكتوب
أعلاه. ومن أجل $f$ ذات القيم في $\R$ تكفي متراجحة القيمة
المتوسطة بمتغير واحد؛ أما في حالة القيم المتجهية فنطبقها على $t \mapsto \langle u,
\varphi(t)\rangle$
حيث $u$ المتجهة الواحدية الموازية للمتجهة $f(b) - f(a)$.
أما الثبوت: فالدالة ثابتة محليا (القطع داخل الكرات) مع الترابط
(المجموعة التي تأخذ فيها $f$ قيمة معطاة \omterm{def:b2:metric:topology}{مفتوحة} ومغلقة في آن:
\cref{ch:b2:metric}).
\end{proof}

\begin{example}[ثابت ليبشيتز انطلاقا من متراجحة القيمة المتوسطة]\label{ex:b2:diffcalc:lipschitz}
هل $f(x, y) = \sin x\,\sin y$ ليبشيتزية على $\R^2$؟ وبأي ثابت؟ تدرجها هو $\nabla f = (\cos x\sin y,\
\sin x\cos y)$، ومربع
معياره
\[
\cos^2x\sin^2y + \sin^2x\cos^2y
\leq \sin^2 y + \cos^2y\cdot 1 = 1
\]
(نحصر $\cos^2x$ و$\sin^2x$ كلا على حدة بالمقدار $1$)، ومنه
$\vertiii{\dd f_{(x,y)}} = \norm{\nabla f} \leq 1$ في كل نقطة، وتطبيق \cref{thm:b2:diffcalc:mvi} على القطعة الواصلة بين نقطتين
كيفما كانتا يعطي
\[
\abs{f(b) - f(a)} \leq \norm{b - a}_2 :
\]
إذن $f$ ليبشيتزية بالثابت $1$، والثابت أمثلي (بجوار المبدأ يبلغ
ميل $f(x, \tfrac\pi2) = \sin x$ القيمة $1$). والخلاصة: متراجحة القيمة المتوسطة تحول
حصرا نقطيا على \omterm{def:b2:diffcalc:differential}{التفاضل} إلى مقياس اتصال شامل --- وهي الطريق
المعتاد إلى تقديرات ليبشيتز في كل بعد، والمحرك الكامن داخل
\cref{exo:b2:diffcalc:12}.
\end{example}

\section{المشتقات الثانية}

\begin{theorem}[شوارتز]\label{thm:b2:diffcalc:schwarz}
إذا كانت $f$ من الصنف $C^2$ على $U$ (أي إن جميع المشتقات
الجزئية الثانية موجودة ومتصلة)، فإن من أجل كل $i, j$:\index{مبرهنة شوارتز}
\[
\frac{\partial^2 f}{\partial x_i\,\partial x_j}
= \frac{\partial^2 f}{\partial x_j\,\partial x_i} .
\]
\end{theorem}

\begin{proof}
يكفي متغيران ($x = x_i$، $y = x_j$، والباقي مجمد).
لننظر في الفرق الثاني
\[
\Delta(h) = f(a + h, b + h) - f(a + h, b) - f(a, b + h) + f(a,b) .
\]
نثبت $h$ ونضع $\varphi(x) = f(x, b+h) - f(x, b)$: عندئذ $\Delta(h)
= \varphi(a + h) - \varphi(a)$، ويعطي تطبيقان لمبرهنة القيمة
المتوسطة
\[
\Delta(h) = h\,\varphi'(\xi)
= h\Bigl(\frac{\partial f}{\partial x}(\xi, b+h) -
\frac{\partial f}{\partial x}(\xi, b)\Bigr)
= h^2\,\frac{\partial^2 f}{\partial y\,\partial x}(\xi, \eta),
\]
مع $\xi \in \intoo{a}{a+h}$، $\eta \in \intoo{b}{b+h}$. وبالاتصال،
$\frac{\Delta(h)}{h^2} \to \frac{\partial^2 f}{\partial
y\partial x}(a, b)$ عندما $h \to 0$. والحساب نفسه بعد تبادل دوري المتغيرين
(بتجميد المتغير الثاني أولا)
يعطي $\frac{\Delta(h)}{h^2} \to \frac{\partial^2 f}{\partial x
\partial y}(a,b)$: فنهايتا المقدار نفسه متساويتان.
\end{proof}

\begin{example}[لماذا نحتاج إلى $C^2$: المثال المضاد لبيانو]\label{ex:b2:diffcalc:peano}
نضع $f(x, y) = \dfrac{xy(x^2 - y^2)}{x^2 + y^2}$، $f(0,0) = 0$.
بعيدا عن المبدأ تكون $f$ من الصنف $C^\infty$؛ وفي المبدأ توجد جميع
المشتقات الجزئية الأولى والثانية، لكن المختلطتين منها متباينتان.
لنحسب على المحورين: $f(x, 0) = f(0, y) = 0$، ومن أجل $y \neq 0$،
\[
\frac{\partial f}{\partial x}(0, y)
= \lim_{x\to0}\frac{f(x,y)}{x}
= \frac{y(0 - y^2)}{y^2} = -y ,
\qquad\text{وبالتناظر}\qquad
\frac{\partial f}{\partial y}(x, 0) = x .
\]
ومنه
\[
\frac{\partial^2 f}{\partial y\,\partial x}(0,0)
= \frac{\dd}{\dd y}\Bigl[\frac{\partial f}{\partial
x}(0,y)\Bigr]_{y=0} = -1,
\qquad
\frac{\partial^2 f}{\partial x\,\partial y}(0,0) = +1 :
\]
فالمشتقتان المختلطتان موجودتان ومختلفتان. ولا تناقض مع
\cref{thm:b2:diffcalc:schwarz}: فالمشتقات الجزئية الثانية للدالة $f$ غير متصلة
في $0$ (اختبرها على $y = tx$). والخلاصة: مبرهنة شوارتز مبرهنة
حقيقية عن \emph{الاتصال}، لا متطابقة شكلية --- والفرضية ``$C^2$''
في مبرهنة تايلور--يونغ الآتية تؤدي عملا حقيقيا.
\end{example}

\begin{theorem}[تايلور--يونغ من الرتبة 2]\label{thm:b2:diffcalc:taylor}
لتكن $f \colon U \to \R$ من الصنف $C^2$ ولتكن $a \in U$. عندئذ، عندما $h \to 0$،
\[
f(a + h) = f(a) + \langle\nabla f(a), h\rangle
+ \frac12\, \langle H_a h,\, h\rangle + o\bigl(\norm h^2\bigr),
\]
حيث $H_a = \bigl(\frac{\partial^2 f}{\partial x_i\partial
x_j}(a)\bigr)$ هي \emph{مصفوفة هس}\index{مصفوفة هس} (وهي متناظرة، بحسب
شوارتز).
\end{theorem}

\begin{proof}
نطبق مبرهنة تايلور--يونغ بمتغير واحد (مجلد السنة الأولى) على
$\varphi(t) = f(a + th)$ على $\intcc{0}{1}$: بقاعدة السلسلة،
$\varphi'(t) = \langle \nabla f(a + th), h\rangle$ و$\varphi''(t)
= \langle H_{a+th}h, h\rangle$، وكلاهما \omterm{def:b2:metric:continuity}{متصل} بدلالة $t$. ثم
$\varphi(1) = \varphi(0) + \varphi'(0) + \frac12\varphi''(\theta)$
(تايلور--لاغرانج) مع $\theta \in \intoo{0}{1}$، ويحول اتصال المشتقات الجزئية الثانية
المقدار $\varphi''(\theta) =
\varphi''(0) + o(1)\cdot\norm h^2$ بانتظام: وهو النشر المكتوب أعلاه.
\end{proof}

\begin{example}[نشر بطريقتين]\label{ex:b2:diffcalc:taylortwoways}
لننشر $f(x, y) = \eu^x\cos y$ في المبدأ حتى الرتبة $2$.
\emph{بتركيب النشور ذات المتغير الواحد:}
\[
\eu^x\cos y = \Bigl(1 + x + \frac{x^2}{2} +
o(x^2)\Bigr)\Bigl(1 - \frac{y^2}{2} + o(y^2)\Bigr)
= 1 + x + \frac{x^2 - y^2}{2} + o\bigl(\norm{(x,y)}^2\bigr) .
\]
\emph{بالمشتقات الجزئية:} $f_x = \eu^x\cos y$، $f_y =
-\eu^x\sin y$، ومنه $\nabla f(0) = (1, 0)$؛ ثم $f_{xx} = f$،
$f_{yy} = -f$، $f_{xy} = -\eu^x\sin y$ تعطي $H_0 =
\operatorname{diag}(1, -1)$: فالمقدار \cref{thm:b2:diffcalc:taylor}
يعيد إنتاج $1 + x + \frac12(x^2 - y^2)$. والحسابان متوافقان، وطريق التركيب كان أسرع ---
دون أي مشتقة جزئية ثانية. والخلاصة: المبدأ \emph{ليس} نقطة حرجة
($\nabla f \neq 0$)، فرغم أن مصفوفة هس غير معرفة الإشارة لا يوجد سرج نعلنه:
الحد الخطي هو المتحكم، والاختبار من الرتبة الثانية لا ينطق إلا في
النقاط الحرجة.
\end{example}

\begin{theorem}[اختبار القيم القصوى من الرتبة الثانية، مبرهنا]\label{thm:b2:diffcalc:extrema}
لتكن $f$ من الصنف $C^2$ بجوار نقطة حرجة $a$ ($\nabla f(a) = 0$)،
ولتكن مصفوفة هس $H = H_a$.
\begin{enumerate}
  \item إذا كانت $H$ معرفة موجبة، فإن $a$ قيمة صغرى محلية تامة
        (ومعرفة سالبة: قيمة عظمى).
  \item إذا كان لدى $H$ قيم ذاتية من الإشارتين، فإن $a$ سرج: لا
        قيمة قصوى.
  \item إذا كانت $H$ شاذة (وشبه معرفة)، فلا نتيجة.
\end{enumerate}
واختبار السنة الأولى ``$rt - s^2$'' هو الحالة $n = 2$: $\det H = rt -
s^2$، وتُقرأ إشارة
$\operatorname{tr}$ من $r$.
\end{theorem}

\begin{proof}
بالمبرهنة الطيفية (\cref{thm:b2:quadratic:spectral})، تكون الصورة التربيعية للمصفوفة $H$
محصورة بين قيمتيها الذاتيتين الحديتين:
$\lambda_{\min}\norm h^2 \leq \langle Hh, h\rangle \leq
\lambda_{\max}\norm h^2$.

(1) إذا كان $\lambda_{\min} > 0$: تعطي مبرهنة تايلور--يونغ
\[
f(a + h) - f(a) \geq \frac{\lambda_{\min}}{2}\norm h^2 -
o(\norm h^2) > 0
\]
من أجل $h \neq 0$ صغيرة: قيمة صغرى محلية تامة.

(2) على \omterm{def:b2:reduction:eigen}{متجهة ذاتية} $v_+$ مع $\lambda_+ > 0$: لدينا $f(a + tv_+) -
f(a) = \frac{\lambda_+}{2}t^2 + o(t^2) > 0$ من أجل $t$ صغيرة؛
وعلى $v_-$ مع $\lambda_- < 0$ يكون الفرق سالبا: فالإشارتان تظهران معا في كل
جوار.

(3) المقادير $f(x,y) = x^2 + y^4$، $x^2 - y^4$، $x^2 + y^3$ تشترك في مصفوفة هس الشاذة شبه
المعرفة نفسها في $0$، مع ثلاثة سلوكات مختلفة.
\end{proof}

\begin{example}[تصنيف كامل، مع الطابع الشامل]\label{ex:b2:diffcalc:quarticrun}
لنصنف جميع القيم القصوى للدالة $f(x, y) = x^4 + y^4 - 4xy$ على $\R^2$.
\emph{النقاط الحرجة:} $\nabla f = (4x^3 - 4y,\ 4y^3 - 4x) =
0$ يعطي $y = x^3$ و$x = y^3 = x^9$، ومنه $x(x^8 - 1) = 0$:
والحلول الحقيقية هي $(0,0)$، $(1,1)$، $(-1,-1)$.
\emph{مصفوفات هس:} $H = \begin{pmatrix} 12x^2 & -4\\ -4 &
12y^2\end{pmatrix}$. في $(\pm1, \pm1)$: $\begin{pmatrix} 12 &
-4\\ -4 & 12\end{pmatrix}$، وقيمتاها الذاتيتان $8$
و$16$: فهي معرفة موجبة، وهناك قيمتان صغريان محليتان تامتان
بالقيمة $f = -2$. وفي $(0,0)$: $\begin{pmatrix} 0 & -4\\ -4 & 0\end{pmatrix}$، وقيمتاها الذاتيتان $\pm4$: فهي سرج.
\emph{الطابع الشامل:} انطلاقا من $2\abs{xy} \leq x^2 +
y^2$،
\[
f(x, y) \geq x^4 + y^4 - 2(x^2 + y^2)
= (x^2 - 1)^2 + (y^2 - 1)^2 + x^2 + y^2 - 2
\xrightarrow[\norm{(x,y)}\to\infty]{} +\infty :
\]
فالدالة $f$ قسرية، ومنه فهي تبلغ قيمة صغرى شاملة (بتراص مجموعات
تحت المستوى)، وذلك بالضرورة في نقطة حرجة: فالقيمة $-2$، المبلوغة
في $(1,1)$ وفي $(-1,-1)$ معا، هي القيمة الصغرى الشاملة؛ ولا توجد قيمة
عظمى (فالمقدار $f$ غير محدود من الأعلى). والخلاصة: الاختبار
المحلي يصنف المرشحين، لكن حجة نمو وحدها هي التي تحول ``المحلي'' إلى
``الشامل'' --- وهي الخطوتان اللتان يسير عليهما كل برهان أمثلة في هذا
الكتاب.
\end{example}

\begin{method}[تصنيف القيم القصوى للدالة $f \colon \R^n \to \R$]\label{met:b2:diffcalc:extrema}
\begin{enumerate}
  \item نحل $\nabla f = 0$ (فنحصل على جميع النقاط الحرجة؛ وعلى مجال ذي حافة
        نعالج الحافة على حدة كما في
        \cref{exo:b2:diffcalc:7}).
  \item في كل نقطة حرجة، نحسب مصفوفة هس وإشارات قيمها الذاتية ---
        وفي البعد $2$ يكفي $\det H$ و$\operatorname{tr} H$: فإذا كان $\det < 0$ فسرج؛ وإذا
        كان $\det >
        0$ فقيمة قصوى، من النوع الذي تحدده إشارة الأثر؛ وإذا
        كان $\det = 0$ فالاختبار صامت، وندرس $f$ على منحنيات.
  \item من أجل النتائج الشاملة، نضيف حجة تراص أو قسرية ($f \to +\infty$ عند
        اللانهاية، أو مجموعة قيود متراصة)، ثم نقارن القيم الحرجة.
\end{enumerate}
\end{method}

\section{مبرهنة الدالة العكسية}

\begin{theorem}[مبرهنة الدالة العكسية]\label{thm:b2:diffcalc:inverse}
لتكن $f \colon U \to \R^n$ من الصنف $C^1$ ولتكن $a \in U$ بحيث يكون $\dd f_a$
\emph{قابلا للقلب}. عندئذ يوجد جواران مفتوحان $V \ni a$، $W
\ni f(a)$ بحيث
يكون $f \colon V \to W$ تقابلا عكسه من الصنف $C^1$، ولدينا\index{مبرهنة الدالة العكسية}
\[
\dd (f^{-1})_{f(x)} = (\dd f_x)^{-1} \qquad (x \in V).
\]
\end{theorem}
\admitted

\begin{remark}[لماذا هي صحيحة: استراتيجية النقطة الثابتة]
حل $f(x) = y$ بجوار $a$ يُعاد كتابته على صورة معادلة النقطة الثابتة
$x = x + \dd f_a^{-1}\bigl(y - f(x)\bigr) =: \Phi_y(x)$؛ والتطبيق
$\Phi_y$ تفاضله $\mathrm{id} - \dd f_a^{-1}\dd f_x$، وهو صغير بجوار $a$ باتصال $\dd f$، ومنه فإن
$\Phi_y$ تطبيق تقلصي على كرة مغلقة صغيرة، فتوفر مبرهنة نقطة بناخ
الثابتة (\cref{thm:b2:metric:banach}) الحل المحلي الوحيد $x = f^{-1}(y)$. أما اتصال الدالة
العكسية وقابليتها \omterm{def:b2:diffcalc:differential}{للتفاضل} فينتجان من التقديرات الواردة في التقلص.
والحساب التفصيلي الكامل يُنجز في السنة الثالثة؛ أما الاستراتيجية ---
والنص --- فتُستعملان بحرية من الآن فصاعدا. وتنتج المبرهنة الرفيقة،
\emph{مبرهنة الدالة الضمنية} (وهي حل $F(x, y) = 0$ بدلالة $y(x)$ عندما
يكون $\frac{\partial F}{\partial y}$ قابلا للقلب)، بتطبيق المبرهنة على $(x, y) \mapsto (x, F(x,y))$.
\end{remark}

\begin{example}[الإحداثيات القطبية]\label{ex:b2:diffcalc:polar}
\omterm{def:b2:diffcalc:differential}{مصفوفة جاكوبي} للتطبيق $\Phi(r, \theta) = (r\cos\theta, r\sin\theta)$ هي
\[
J_\Phi = \begin{pmatrix} \cos\theta & -r\sin\theta\\ \sin\theta &
r\cos\theta \end{pmatrix},
\qquad \det J_\Phi = r :
\]
وهي قابلة للقلب من أجل $r \neq 0$، ومنه فإن $\Phi$ تقابل تفاضلي محلي من
الصنف $C^1$ بعيدا عن المبدأ --- وهو الترخيص الذي يبيح ``الانتقال
إلى الإحداثيات القطبية''، مجددا من أجل التكاملات المتعددة في
\cref{ch:b2:multint}.
\end{example}

\begin{example}[محلي في كل مكان، شامل في لا مكان]\label{ex:b2:diffcalc:localnotglobal}
نضع $f(x, y) = \bigl(\eu^x\cos y,\ \eu^x\sin y\bigr)$ على
$\R^2$. \omterm{def:b2:diffcalc:differential}{مصفوفة جاكوبي}،
\[
J_f = \begin{pmatrix}
\eu^x\cos y & -\eu^x\sin y\\
\eu^x\sin y & \eu^x\cos y
\end{pmatrix},
\qquad
\det J_f = \eu^{2x} > 0 ,
\]
لا تنعدم أبدا: فبحسب \cref{thm:b2:diffcalc:inverse} يكون $f$ تقابلا تفاضليا محليا من
الصنف $C^1$ في \emph{كل} نقطة من المستوي.
ومع ذلك فإن $f$ بعيد كل البعد عن التباين: لدينا $f(x, y + 2\pi) = f(x, y)$، فكل قيمة
تُبلغ عددا لانهائيا من المرات؛ وهو ليس غامرا أيضا، إذ إن $\norm{f(x, y)} = \eu^x > 0$
يفوّت المبدأ. والخلاصة: مبرهنة الدالة العكسية \emph{محلية} على نحو
لا يُختزل --- فقابلية قلب كل $\dd f_a$ تعطي ترقيعا من دوال عكسية محلية
ليس من الضروري أن تتجمع في واحدة.
(وسيتعرف القراء الذين يعرفون الأعداد المركبة على $z \mapsto
\eu^z$؛ والترقيع هو
عائلة فروع اللوغاريتم.)
قارن ذلك بالتمرين \cref{exo:b2:diffcalc:12}، حيث تفرض فرضية شاملة كمية وجود دالة عكسية
شاملة واحدة.
\end{example}

\begin{remark}[مزالق شائعة]
\emph{(1) المشتقات الاتجاهية رخيصة، والتفاضلات ليست كذلك:} قد توجد
جميع المشتقات الاتجاهية --- بل وقد لا تتعلق خطيا بالاتجاه --- دون
قابلية \omterm{def:b2:diffcalc:differential}{للتفاضل} (\cref{exo:b2:diffcalc:2})؛ ووحده محك $C^1$
(\cref{thm:b2:diffcalc:c1}) يرفع المشتقات الجزئية إلى \omterm{def:b2:diffcalc:differential}{تفاضل}.
\emph{(2) الحرج لا يعني الأقصى:} فالسروج (\cref{ex:b2:diffcalc:quarticrun}) والحالة الشاذة
الصامتة (\cref{thm:b2:diffcalc:extrema}~(3)) كلتاهما تختبئان
وراء $\nabla f = 0$. \emph{(3) لا مساواة قيمة متوسطة ذات قيم متجهية:}
وحدها \emph{المتراجحة} في
\cref{thm:b2:diffcalc:mvi} تبقى صحيحة
(\cref{exo:b2:diffcalc:9})؛ ولا تكتب أبدا $f(b) - f(a) = \dd
f_c(b-a)$ من أجل $f$ ذات قيم في $\R^m$، $m \geq 2$.
\emph{(4) قابلية القلب المحلية ليست تباينا:}
\cref{ex:b2:diffcalc:localnotglobal}. \emph{(5) التدرج ملك للجداء السلمي:} فالمقدار $\nabla f$ هو
المتجهة الممثلة \omterm{def:b2:diffcalc:differential}{للتفاضل} $\dd f_a$ في جداء سلمي مختار؛ فإن غيّرنا الجداء
(كما في المثال الموزون من فصل \omterm{def:b2:quadratic:def}{الصور التربيعية}) دار التدرج، بينما
\omterm{def:b2:diffcalc:differential}{التفاضل} --- وهو الكائن الجوهري --- لا يتحرك.
\end{remark}

\begin{remark}[أين يُستعمل هذا]
كل ما يأتي بعد هذا الفصل هو حساب \omterm{def:b2:diffcalc:differential}{تفاضل} مطبق: ففصل المعادلات
التفاضلية يخطّط التدفقات ويستعمل صيغة ليوفيل \omterm{def:b2:linalg:det}{للمحدد} (المبرهنة في
مسألة نهاية الأسبوع لهذا الفصل)؛ وفصلا المنحنيات والسطوح يدرسان
مجموعات المستوى والتوسيمات عبر مبرهنة الدالة الضمنية؛ والتكاملات
المتعددة تغير المتغيرات بمصفوفات جاكوبي.
وتطور مسألة نهاية الأسبوع حساب \omterm{def:b2:diffcalc:differential}{التفاضل} \emph{على فضاء المصفوفات
نفسه} --- \omterm{def:b2:diffcalc:differential}{تفاضل} \omterm{def:b2:linalg:det}{المحدد}، \omterm{def:b2:diffcalc:differential}{وتفاضل} المقلوب، والأسي المصفوفي، والزمرة
المتعامدة بوصفها مجموعة مستوى ملساء --- وهو ظل السنة الثانية لما
يصوغه مجلد السنة الثالثة على أنه متنوعات وزمر لي.
\end{remark}

\section{تمارين}

\begin{exercise}[$\star$]\label{exo:b2:diffcalc:1}
احسب مصفوفتي جاكوبي للتطبيقين $f(x,y) = (x^2 - y^2,\, 2xy)$ و$\Phi(r,\theta,z) = (r\cos\theta, r\sin\theta, z)$؛ وأين يكون التفاضلان
قابلين للقلب؟
\end{exercise}

\begin{exercise}[$\star$]\label{exo:b2:diffcalc:2}
لتكن $f(x,y) = \frac{x^3}{x^2 + y^2}$ ($f(0,0) = 0$). برهن على أن جميع المشتقات الاتجاهية للدالة
$f$ في $0$ موجودة، لكن $f$ ليست \omterm{def:b2:diffcalc:differential}{قابلة للتفاضل} في $0$
\emph{(التطبيق $v \mapsto$ المشتقة الاتجاهية ليس خطيا)}.
\end{exercise}

\begin{exercise}[$\star$]\label{exo:b2:diffcalc:3}
جد النقاط الحرجة للدالة $f(x, y) = x^3 + y^3 -
3xy$ وصنفها باستعمال \cref{thm:b2:diffcalc:extrema}، ثم افعل
الشيء نفسه من أجل $g(x,y) = x^4 +
y^4 - 2(x - y)^2$.
\end{exercise}

\begin{exercise}[$\star\star$]\label{exo:b2:diffcalc:4}
لتكن $f \colon \R^n \to \R$ من الصنف $C^1$ و\emph{متجانسة من الدرجة $p$}:
$f(tx) = t^pf(x)$ من أجل $t > 0$. برهن على متطابقة أويلر
\[
\langle \nabla f(x), x\rangle = p\,f(x) ,
\]
وعلى عكسها من أجل الدوال من الصنف $C^1$ على $\R^n\setminus\{0\}$.
\end{exercise}

\begin{exercise}[$\star\star$]\label{exo:b2:diffcalc:5}
لتكن $A$ متناظرة ولتكن $f(x) = \frac12\langle Ax, x\rangle -
\langle b, x\rangle$. احسب $\nabla f$ و$H_f$؛ ومتى تكون
$f$ محدبة؟ بافتراض $A$ معرفة موجبة، بيّن أن $f$ تقبل قيمة
صغرى شاملة وحيدة في حل المعادلة $Ax = b$ --- وهي علة وجود خوارزمية
النزول بالتدرج.
\end{exercise}

\begin{exercise}[$\star\star$]\label{exo:b2:diffcalc:6}
(مضاعف لاغرانج، بقيد واحد، مبرهنا يدويا) لتكن $f, g$ من الصنف
$C^1$ على $\R^2$، ولنفترض أن $f$ تبلغ في $a$ قيمة قصوى محلية
على مجموعة المستوى $\{g = 0\}$، مع $\nabla g(a) \neq 0$. برهن على أن $\nabla f(a) = \lambda\nabla g(a)$ من أجل
$\lambda$ ما.
\emph{(وسّم مجموعة المستوى بجوار $a$ بمبرهنة الدالة الضمنية ثم
فاضل $t \mapsto f(\gamma(t))$.)} تطبيق: القيم القصوى للدالة $f(x,y) = xy$ على الدائرة
$x^2 + y^2 = 1$.
\end{exercise}

\begin{exercise}[$\star\star$]\label{exo:b2:diffcalc:7}
عيّن القيم القصوى للدالة $f(x, y) = x^2 + y^2 - xy + x - y$ على $\R^2$، ثم قيمتيها العظمى
والصغرى على المثلث المغلق ذي الرؤوس $(0,0)$ و$(1,0)$ و$(0,1)$
\emph{(النقاط الحرجة الداخلية، ثم الأضلاع الثلاثة، ثم الرؤوس)}.
\end{exercise}

\begin{exercise}[$\star\star\star$]\label{exo:b2:diffcalc:8}
لتكن $f \colon \R^2 \to \R^2$، $f(x, y) = (x + y^2,\; y + x^2)$.
بيّن أن $f$ تقابل تفاضلي محلي بجوار $0$، واحسب $\dd(f^{-1})_{(0,0)}$، وجد
أكبر $r$ بحيث يكون $\dd f$ قابلا للقلب على الكرة $\norm{(x,y)}_2 < r$
\emph{(احسب $\det J_f$)}.
\end{exercise}

\begin{exercise}[$\star\star\star$]\label{exo:b2:diffcalc:9}
(مبرهنة رول تسقط، ومبرهنة القيمة المتوسطة تصمد) أعط $f \colon \R \to \R^2$، $C^1$،
مع $f(0) = f(2\pi)$ لكن $f'(t) \neq 0$ من أجل كل $t$ (فلا مبرهنة رول ذات قيم
متجهية). ثم تحقق على مثالك من \emph{متراجحة} القيمة المتوسطة في
\cref{thm:b2:diffcalc:mvi}.
\end{exercise}

\begin{exercise}[$\star$]\label{exo:b2:diffcalc:10}
احسب \omterm{def:b2:diffcalc:differential}{التفاضل} والتدرج للدالتين $f(x) = \norm
x_2^2$ و$g(x) = \langle Ax, x\rangle$ على $\R^n$ (حيث $A$
مصفوفة مربعة لا نفترضها متناظرة)، ثم احسب مصفوفة هس لكل منهما.
ومن أجل أي $A$ تكون $g$ محدبة؟
\end{exercise}

\begin{exercise}[$\star\star$]\label{exo:b2:diffcalc:11}
لتكن $f \colon \R^n \to \R$ من الصنف $C^2$. برهن على أن $f$ محدبة إذا وفقط إذا
كانت مصفوفة هس $H_x$ شبه معرفة موجبة في كل $x$
\emph{(اختزل الأمر إلى متغير واحد: $t \mapsto f(a + t(b-a))$؛ واستعمل تايلور--لاغرانج
في اتجاه واحد، وقيّم $\varphi''$ من أجل العكس)}.
\end{exercise}

\begin{exercise}[$\star\star\star$]\label{exo:b2:diffcalc:12}
(مبرهنة قلب شاملة) لتكن $g \colon \R^n \to \R^n$ من الصنف $C^1$
مع $\vertiii{\dd g_x} \leq k < 1$ من أجل كل $x$، ولتكن $f =
\mathrm{id} + g$.
\begin{enumerate}
  \item بيّن أن $\norm{f(x) - f(y)} \geq (1 - k)\norm{x - y}$: ومنه فإن $f$ متباين، وعكسه \omterm{def:b2:metric:continuity}{متصل} على صورته.
  \item بيّن أن التطبيق $x \mapsto y -
        g(x)$ تقلصي على الفضاء التام $\R^n$ من أجل
        كل $y \in \R^n$، واستنتج بمبرهنة نقطة بناخ الثابتة
        (\cref{thm:b2:metric:banach}) أن $f$ غامر.
  \item استنتج أن $f$ تقابل من $\R^n$ على نفسه، عكسه \omterm{def:b2:metric:continuity}{ليبشيتزي}
        بالثابت $(1-k)^{-1}$ --- وهو النظير \emph{الشامل} للمبرهنة
        \cref{thm:b2:diffcalc:inverse} (التي هي، على النقيض، محلية بحتة).
\end{enumerate}
\end{exercise}

\section{مسألة: حساب تفاضل المصفوفات --- جاكوبي والأسي والزمرة
المتعامدة}

\begin{problem}\label{pb:b2:diffcalc:1}
إن أنقى ميدان للعب حساب \omterm{def:b2:diffcalc:differential}{التفاضل} هو الفضاء $\mathcal M_n(\R) \simeq \R^{n^2}$ نفسه: فأكثر
تطبيقاته طبيعية --- الجداء والمقلوب \omterm{def:b2:linalg:det}{والمحدد} والأسي --- لها تفاضلات
بالغة الأناقة. وتحسب هذه المسألة هذه التفاضلات كلها: \emph{متسلسلة
نويمان}، \omterm{def:b2:diffcalc:differential}{وتفاضل} المقلوب، و\emph{صيغة جاكوبي}\index{صيغة جاكوبي}
\omterm{def:b2:linalg:det}{للمحدد} مع \emph{صيغة ليوفيل} ربحا إضافيا، و\emph{الأسي
المصفوفي}\index{الأسي المصفوفي} مع
$\det\eu^A = \eu^{\operatorname{tr}A}$، وأخيرا الزمرة المتعامدة $O_n$ بوصفها مجموعة مستوى ملساء
مماسها فضاء المصفوفات المتخالفة التناظر --- وهي الهندسة التفاضلية
في مهدها. وفي كل ما يلي، $\vertiii\cdot$ هو \omterm{thm:b2:nvs:continuouslinear}{معيار المؤثر} المرافق \omterm{def:b2:nvs:norm}{للمعيار}
$\norm\cdot_2$، و$\langle X, Y\rangle =
\operatorname{tr}(X^{\mathsf T}Y)$ هو الجداء السلمي لفروبينيوس.

\medskip
\noindent\textbf{الجزء الأول --- متسلسلة نويمان.}
\begin{enumerate}
  \item برهن على الخاصية تحت الضربية $\vertiii{AB} \leq
        \vertiii A\,\vertiii B$، واستنتج أن التطبيقات
        الحدودية بدلالة $A$ (جداءات المصفوفات \omterm{def:b2:linalg:det}{والمحدد} والأثر)
        متصلة على $\mathcal M_n(\R)$.
  \item من أجل $\vertiii X < 1$، بيّن أن $\sum_{k\geq0}X^k$ تتقارب تقاربا مطلقا في $\mathcal M_n(\R)$
        (\cref{thm:b2:nvs:absoluteconvergence})، وأن مجموعها هو $(I - X)^{-1}$، وأن
        \[
        \vertiii{(I - X)^{-1}} \leq \frac{1}{1 - \vertiii X},
        \qquad
        (I - X)^{-1} = I + X + O\bigl(\vertiii X^2\bigr) .
        \]
  \item استنتج أن $GL_n(\R)$ \emph{\omterm{def:b2:metric:topology}{مفتوحة}}: أي إذا كانت $A$ قابلة
        للقلب و$\vertiii H <
        \frac{1}{\vertiii{A^{-1}}}$، فإن $A + H$ قابلة للقلب. واستنتج أيضا أن $GL_n(\R)$
        \emph{كثيفة} في $\mathcal M_n(\R)$ \emph{(اضطرب $A$ بالمقدار $\varepsilon I$:
        فالمقدار $\det(A + \varepsilon I)$ كثير حدود غير معدوم بدلالة $\varepsilon$)}.
  \item بيّن أن تطبيق القلب $\Phi(A) = A^{-1}$ \omterm{def:b2:metric:continuity}{متصل} على $GL_n(\R)$.
\end{enumerate}

\medskip
\noindent\textbf{الجزء الثاني --- التفاضلات الأولى.}
\begin{enumerate}[resume]
  \item بيّن أن تطبيق التربيع $A \mapsto A^2$ قابل \omterm{def:b2:diffcalc:differential}{للتفاضل} وتفاضله $H \mapsto AH + HA$،
        وبوجه أعم أن $A \mapsto A^k$ تفاضله $H \mapsto \sum_{i=0}^{k-1}
        A^iHA^{k-1-i}$. ولماذا \emph{لا يمكن} أن نكتب
        $kA^{k-1}H$ في الحالة العامة؟
  \item برهن على أن $\Phi(A) = A^{-1}$ قابل \omterm{def:b2:diffcalc:differential}{للتفاضل} على $GL_n(\R)$ وأن
        \[
        \dd\Phi_A(H) = -A^{-1}HA^{-1}
        \]
        \emph{(اكتب $(A + H)^{-1} = (I + A^{-1}H)^{-1}A^{-1}$ وانشر بالسؤال 2)}. وتحقق من الصيغة في
        الحالة السلمية $n = 1$.
  \item من أجل منحن $t \mapsto A(t) \in GL_n(\R)$ من الصنف $C^1$، استنتج $\bigl(A(t)^{-1}\bigr)' = -A^{-1}A'A^{-1}$، ثم انشر
        $t \mapsto (I + tB)^{-1}$ حتى الرتبة الأولى في $t = 0$.
  \item احسب \omterm{def:b2:diffcalc:differential}{تفاضل} $f(A) =
        \operatorname{tr}(A^k)$ وعيّن تدرجه من أجل الجداء السلمي
        لفروبينيوس:
        \[
        \dd f_A(H) = k\operatorname{tr}\bigl(A^{k-1}H\bigr),
        \qquad
        \nabla f(A) = k\,\bigl(A^{k-1}\bigr)^{\mathsf T} .
        \]
  \item الأسئلة نفسها من أجل $f(A) = \operatorname{tr}
        (A^{\mathsf T}A) = \norm A_F^2$: \omterm{def:b2:diffcalc:differential}{التفاضل} والتدرج ومصفوفة هس
        (الثابتة)؛ واستنتج أن $\norm\cdot_F^2$ محدبة تماما.
\end{enumerate}

\medskip
\noindent\textbf{الجزء الثالث --- صيغة جاكوبي.}
\begin{enumerate}[resume]
  \item برهن على أن
        \[
        \det(I + H) = 1 + \operatorname{tr}H +
        O\bigl(\vertiii H^2\bigr)
        \]
        \emph{(انشر $\det(e_1 + h_1, \dots, e_n + h_n)$ بالخطية المتعددة بدلالة الأعمدة: فالحدود
        التي تحوي عمودين على الأقل من أعمدة $h$ هي $O(\vertiii H^2)$)}: أي
        $\dd(\det)_I =
        \operatorname{tr}$.
  \item من أجل $A$ قابلة للقلب، استنتج
        \[
        \dd(\det)_A(H) = \det(A)\,
        \operatorname{tr}\bigl(A^{-1}H\bigr) .
        \]
  \item بيّن أن $\frac{\partial\det}{\partial a_{ij}}(A) = C_{ij}$ من أجل \emph{كل} $A$ (قابلة للقلب أو لا)،
        حيث هذا المقدار هو المرافق الجبري $(i,j)$ \emph{(نشر لابلاس
        وفق السطر $i$)}، ومنه، بدلالة المصفوفة الملحقة $\operatorname{adj}A =
        \operatorname{com}(A)^{\mathsf T}$:
        \[
        \dd(\det)_A(H) =
        \operatorname{tr}\bigl(\operatorname{adj}(A)\,H\bigr),
        \qquad
        \nabla(\det)(A) = \operatorname{com}(A) ,
        \]
        وهو ما يستعيد السؤال 11 عندما تكون $A$ قابلة للقلب
        ($\operatorname{adj}A = \det(A)A^{-1}$). وهذه هي \emph{صيغة جاكوبي}: $\bigl(\det
        A(t)\bigr)' = \operatorname{tr}\bigl(
        \operatorname{adj}(A(t))\,A'(t)\bigr)$.
  \item (صيغة ليوفيل) ليكن $A(t)$ منحنى مصفوفات من الصنف $C^1$
        يحقق المعادلة التفاضلية الخطية
        $A'(t) = M(t)A(t)$. برهن على أن
        \[
        \bigl(\det A(t)\bigr)' =
        \operatorname{tr}\bigl(M(t)\bigr)\,\det A(t),
        \qquad\text{ومنه}\qquad
        \det A(t) = \det A(0)\,
        \exp\Bigl(\int_0^t\operatorname{tr}M\Bigr)
        \]
        \emph{(استعمل $\operatorname{adj}(A)\,A = \det(A)I$ وثبات الأثر بالدوران الدائري)} --- وهي
        متطابقة الفرونسكيان التي سيستعملها فصل المعادلات التفاضلية
        باتصال.
  \item بيّن أن $SL_n(\R) = \{\det = 1\}$ مجموعة مستوى ملساء: أي إن \omterm{def:b2:diffcalc:differential}{التفاضل} $\dd(\det)_A$ في
        كل $A \in SL_n(\R)$ تطبيق خطي \emph{غامر} على $\R$
        \emph{(قيّمه في $H = \frac1nA$)}.
\end{enumerate}

\medskip
\noindent\textbf{الجزء الرابع --- الأسي المصفوفي.}
\begin{enumerate}[resume]
  \item بيّن أن $\eu^A = \sum_{k\geq0}\frac{A^k}{k!}$ تتقارب تقاربا مطلقا من أجل كل $A$، وتقاربا
        ناظميا على كل كرة، مع $\vertiii{\eu^A} \leq
        \eu^{\vertiii A}$؛ وأن $\eu^A$ يتعلق باتصال
        بالمتغير $A$.
  \item برهن على أن $AB = BA$ يستلزم $\eu^{A+B} =
        \eu^A\eu^B$ \emph{(\omterm{thm:b2:series:fubini}{جداء كوشي}، وهو مشروع
        بالتقارب المطلق)}؛ واستنتج أن $\eu^A$ قابل للقلب دائما، وأن
        مقلوبه $\eu^{-A}$: أي إن $\exp$ يرسل $\mathcal M_n(\R)$ داخل $GL_n(\R)$.
  \item بيّن أن $t \mapsto \eu^{tA}$ من الصنف $C^1$ (بل $C^\infty$) وأن
        \[
        \frac{\dd}{\dd t}\,\eu^{tA} = A\,\eu^{tA} =
        \eu^{tA}A
        \]
        \emph{(فاضل المتسلسلة حدا حدا على القطع: فالمتسلسلة المشتقة
        تتقارب تقاربا ناظميا)}.
  \item برهن على المتطابقة
        \[
        \det\bigl(\eu^{A}\bigr) = \eu^{\operatorname{tr}A}
        \]
        \emph{(طبق صيغة ليوفيل، السؤال 13، على $A(t) = \eu^{tA}$)}. وللتحقق:
        $n = 1$؛ ومن أجل $A$ معدومة القوى؛ ثم إن المصفوفات معدومة
        الأثر ترسو في $SL_n(\R)$.
  \item بيّن أن $\eu^H = I + H + O(\vertiii H^2)$، ومنه فإن $\exp$ قابل \omterm{def:b2:diffcalc:differential}{للتفاضل} في $0$ مع
        $\dd(\exp)_0 =
        \mathrm{id}$؛ واستنتج بمبرهنة الدالة العكسية (\cref{thm:b2:diffcalc:inverse}) أن $\exp$
        تقابل تفاضلي من الصنف $C^1$ من جوار للمصفوفة $0$ على
        جوار للمصفوفة $I$: أي إن كل مصفوفة قريبة من مصفوفة الوحدة
        تقبل لوغاريتما.
  \item بيّن أن $\exp$ يرسل المصفوفات المتناظرة على المصفوفات
        المتناظرة \emph{المعرفة الموجبة}، وذلك تقابليا \emph{(قطّر؛
        والتطبيق العكسي هو اللوغاريتم الطيفي)}.
\end{enumerate}

\medskip
\noindent\textbf{الجزء الخامس --- الزمرة المتعامدة بوصفها مجموعة
مستوى.}
\begin{enumerate}[resume]
  \item نضع $F(A) = A^{\mathsf T}A$، من $\mathcal M_n(\R)$ نحو فضاء المصفوفات المتناظرة $S_n$.
        احسب $\dd F_A(H) =
        A^{\mathsf T}H + H^{\mathsf T}A$ وبيّن أن هذا \omterm{def:b2:diffcalc:differential}{التفاضل} \emph{غامر} على $S_n$ في كل
        $A \in O_n = F^{-1}(I)$ \emph{(من أجل $S \in S_n$ معطاة، جرّب $H = \frac12 AS$)}: ومنه فإن
        $O_n$ مجموعة مستوى ملساء بعدها $n^2 - \frac{n(n+1)}2 = \frac{n(n-1)}2$.
  \item بيّن أن كل منحن $A(t) \in O_n$ من الصنف $C^1$ يحقق $A(0) =
        I$ له سرعة
        \emph{متخالفة التناظر} $A'(0)$، وبالعكس أن المنحنى $\eu^{tK}$ يبقى
        في $O_n$ من أجل $K$ متخالفة التناظر: أي إن الفضاء المماس
        للمجموعة $O_n$ في $I$ هو بالضبط فضاء المصفوفات المتخالفة
        التناظر.
  \item بيّن أن $\det\eu^{K} = 1$ من أجل $K$ متخالفة التناظر (السؤال 18):
        فالمنحنى الأسي يعيش في زمرة الدورانات $SO_n$. واحسبه كاملا من
        أجل $n = 2$: بوضع $J =
        \begin{pmatrix}0 & -1\\ 1 & 0\end{pmatrix}$، برهن على أن
        \[
        \eu^{\theta J} = \begin{pmatrix} \cos\theta &
        -\sin\theta\\ \sin\theta & \cos\theta\end{pmatrix} :
        \]
        أي إن الأسي المصفوفي \emph{هو} الدوران بالزاوية $\theta$، وأن
        تعريفي جيب التمام والجيب بالمتسلسلات يظهران من جديد داخل
        مصفوفة.
  \item (حيلة الكثافة) باستعمال كثافة $GL_n(\R)$ (السؤال 3) والاتصال،
        عمّم المتطابقة
        \[
        \operatorname{adj}(AB) =
        \operatorname{adj}(B)\operatorname{adj}(A)
        \]
        من المصفوفات القابلة للقلب إلى جميع المصفوفات \emph{(من أجل
        $A, B$ قابلة للقلب يتساوى الطرفان بالمقدار $\det(AB)(AB)^{-1}$؛ والطرفان
        كلاهما كثير حدود بدلالة المعاملات)}.
  \item تركيب. في جملة واحدة عن كل بند: (1) أي الفصول السابقة وفّر
        محرك كل جزء (التمام والجبور المعيارية؛ المبرهنة الطيفية؛
        مبرهنة الدالة العكسية)؛ (2) أي صيغة من هذه المسألة سيستند
        إليها فصل المعادلات التفاضلية، وأين؛ (3) ماذا يقول $\dd(\det)_I = \operatorname{tr}$
        و$\det\eu^A =
        \eu^{\operatorname{tr}A}$ عن الأثر \omterm{def:b2:linalg:det}{والمحدد} بوصفهما ``الحجم اللامتناهي الصغر
        والحجم الشامل''؛ (4) ماذا يصنع مجلد السنة الثالثة بالأسئلة
        21--23 (زمر لي وجبورها).
\end{enumerate}
\end{problem}
